% !TEX root = ../main.tex
% §6 Experimental Setup. Search: results/stage6_main/search_log.txt (config ids follow run.py GRID order);
% environment.csv and latency.csv, peak memory from main_log.txt. The main run (stages 5-8) used the
% environment in environment.csv (see README); requirements.txt pins the stage 1-4 environment.
\section{Experimental Setup}
\label{sec:setup}

\paragraph{Split and leakage prevention.}
All tiers use the single split of \S\ref{sec:data} (fixed indices, seed~0): 7,705 / 1,194 / 1,682
images for train / val / test.
Models are fitted on train normals only.
Preprocessing uses fixed ImageNet statistics, so nothing is fitted on val or test.
Val anomalies are used only to choose $\tau$ and the hyperparameters below, never to fit weights.
For every object, all models are fitted and all thresholds fixed before the test features are loaded,
and test is evaluated once.
The only earlier contact with test was a pipeline smoke test on one object (candle, seed~0, default
settings) before the search; nothing was chosen from it.

\paragraph{Hyperparameters.}
The search uses val only (seed~0).
For each tier it fits every point of the grid in Table~\ref{tab:hparams} and keeps the setting with the
highest mean val PR-AUC over the 12 objects.
The chosen Tier~2 setting lies at the edge of its grid; we kept the grid as defined in advance.

\begin{table}[t]
  \centering
  \small
  \setlength{\tabcolsep}{4pt}
  \begin{tabular}{@{}llll@{}}
    \toprule
    Tier & Hyperparameter & Range & Selected \\
    \midrule
    T2 & channels $d$ & 50, 100, 200 & 200 \\
    T2 & layers & \{2,3\}, \{2,3,4\} & \{2,3,4\} \\
    T2 & ridge $\epsilon$ & fixed & 0.01 \\
    \midrule
    T3 & coreset ratio & 1\%, 10\% & 1\% \\
    T3 & layers & \{2,3\}, \{2,3,4\} & \{2,3\} \\
    T3 & neighbourhood / proj. & fixed & $3 \times 3$ / 128-d \\
    \bottomrule
  \end{tabular}
  \caption{Hyperparameter search on val (seed~0, full grid per tier). Mean val PR-AUC of the selected
    setting: 0.664 (T2) and 0.811 (T3).}
  \label{tab:hparams}
\end{table}

\paragraph{Metrics.}
PR-AUC is the primary metric because, unlike ROC-AUC, it reflects how many alerts are real defects under
class imbalance \citep{davis2006relationship,saito2015precision}.
We also report ROC-AUC for comparison with the literature.
At the chosen $\tau$ we report F1, F2 and the false-positive rate, which is the share of good products
sent to re-inspection.
We also report $\mathrm{prec}_\pi$ (Eq.~\ref{eq:prec-pi}) at $\pi = 1\%$ and $5\%$, the share of alerts
that would be real defects in production.
Every number is the mean $\pm$ std over seeds 0, 1 and 2 of the 12-object average.
The seed controls the Tier~2 channel subset and the Tier~3 coreset start and projection; Tier~1 is
deterministic.

\paragraph{Implementation and hardware.}
PyTorch 2.15.0 with CUDA 13.2, torchvision 0.29.2 and scikit-learn, on an NVIDIA RTX 4060 Ti (8\,GB).
Latency is the median over 50 val images of pcb1, on GPU with batch size~1 after 10 warm-up images,
including feature extraction.
Training time covers model fitting only, not loading the cached features.
The main run peaked at 2.24\,GB of GPU memory and 6.6\,GB of RAM.
